\documentclass[11pt]{article}
\usepackage[letterpaper,margin=0.82in]{geometry}
\usepackage{amsmath,amssymb,amsthm}
\usepackage{graphicx}
\usepackage[hidelinks]{hyperref}
\usepackage{microtype}
\usepackage[T1]{fontenc}

\newtheorem{theorem}{Theorem}
\newcommand{\Mult}{\operatorname{Mult}}

\title{A Hurwitz counterexample to cyclic Kaplansky approximants}
\author{Candidate answer to Question 9.1 of arXiv:2105.01110}
\date{13 August 2026}

\begin{document}
\maketitle

\paragraph{Status.}
\textbf{Candidate counterexample; likely valid, pending expert review.}
Question 9.1 of Aleman--Hartz--McCarthy--Richter~\cite{AHMR} has a negative
answer as literally stated, already for the Hardy space on the disk.

\section{Source question}

For a normalized complete Pick space $\mathcal H$, write $[m]$ for the
$\mathcal H$-closure of $\Mult(\mathcal H)m$.  Lemma 3.3 of the source says
that if
\[
 m\in\Mult(\mathcal H),\qquad
 \varphi\in\Mult(\mathcal H)\cap[m],\qquad
 \|\varphi\|_{\Mult(\mathcal H)}\leq1,
\]
then there are $\varphi_n\in\Mult(\mathcal H)$ such that
\[
 \varphi_nm\longrightarrow\varphi\quad\hbox{weak-*},
 \qquad \|\varphi_nm\|_{\Mult(\mathcal H)}\leq1.
\]
Question 9.1 asks whether the $\varphi_n$ can always be chosen cyclic
\cite[Question 9.1]{AHMR}.

\begin{center}
\includegraphics[width=0.98\textwidth]{figures/open_problem_crop.png}
\end{center}
\noindent\emph{Source evidence:} arXiv:2105.01110, PDF p.~25,
Question 9.1.

\section{Proof intuition}

Set $m=1$.  The proposed strengthening would then make every contractive
multiplier a bounded weak-* limit of cyclic multipliers.  But cyclic
multipliers cannot vanish: all their multiplier multiples share any interior
zero, so they could not approximate the constant function $1$ in the Hilbert
space.  Bounded weak-* convergence of multipliers implies locally uniform
convergence.  The target $\varphi(z)=z$ has a zero without being identically
zero, and Hurwitz's theorem says it cannot be a locally uniform limit of
zero-free holomorphic functions.

\section{Counterexample}

\begin{theorem}\label{thm:counter}
The improvement proposed in Question 9.1 is false.  One may take
\[
 \mathcal H=H^2(\mathbb D),\qquad m=1,\qquad\varphi(z)=z.
\]
\end{theorem}

\begin{proof}
The Hardy space $H^2(\mathbb D)$ is a normalized complete Pick space, with
kernel $(1-z\overline w)^{-1}$, and its multiplier algebra is
$H^\infty(\mathbb D)$.  The multiplier $m=1$ is cyclic.  Moreover,
$\varphi(z)=z$ belongs to $H^\infty\cap[1]$ (indeed, it is itself a
multiplier multiple of $1$), and $\|\varphi\|_\infty=1$.  Thus all hypotheses
of Lemma 3.3 are satisfied.

Suppose its conclusion could be strengthened as in Question 9.1.  There
would be cyclic multipliers $\varphi_n\in H^\infty$ with
\begin{equation}\label{eq:conv}
 \|\varphi_n\|_\infty\leq1,
 \qquad \varphi_n\longrightarrow z\quad\text{weak-* in }H^\infty.
\end{equation}
Each $\varphi_n$ is zero-free on $\mathbb D$.  Indeed, if
$\varphi_n(a)=0$ for some $a\in\mathbb D$, then every element of
$H^\infty\varphi_n$ vanishes at $a$.  Point evaluation is continuous on
$H^2$, so the $H^2$-closure $[\varphi_n]$ also consists of functions
vanishing at $a$ and cannot contain $1$.  This contradicts cyclicity.

On bounded subsets of the multiplier algebra, weak-* convergence implies
pointwise convergence; this fact is also recalled immediately before Lemma
3.3 in the source~\cite{AHMR}.  The uniform bound in~\eqref{eq:conv} and
Montel's theorem then upgrade the convergence to local uniform convergence.
For completeness, every subsequence has a locally uniformly convergent
subsequence, and its pointwise limit must be $z$; this forces the whole
sequence to converge locally uniformly to $z$.

Hurwitz's theorem~\cite[Theorem 10.28]{Rudin} now says that the locally
uniform limit of the zero-free holomorphic functions $\varphi_n$ is either
zero-free or identically zero.  The function $z$ is neither: it vanishes at
$0$ and is not the zero function.  This contradiction proves the claim.
\end{proof}

\paragraph{Stronger obstruction.}
The same proof shows that, with $m=1$, no nonzero contractive multiplier
having an interior zero can be the limit required by the proposed cyclic
strengthening.

\section{Verification, limitations, and novelty}

\begin{itemize}
\item Every hypothesis is literal: $H^2(\mathbb D)$ is a normalized complete
Pick space; $m=1$ and $\varphi=z$ are multipliers; $m$ is cyclic;
$\varphi\in[m]$; and $\|\varphi\|=1$.
\item No identification of weak-* and local-uniform topologies is assumed
without boundedness.  The norm bound is part of Lemma 3.3, and Montel is used
only after that bound.
\item The counterexample answers the formal Question 9.1, which asks to
strengthen Lemma 3.3 for arbitrary admissible $\varphi$.  It does \emph{not}
settle the narrower preceding prose question: for an arbitrary cyclic $m$,
can cyclic $\varphi_n$ be found with $\varphi_nm\to1$?  Our target is
$\varphi=z$, not $1$.
\item The parsed source is revised 21 April 2026 and still states Question
9.1.  Bounded run-index and web/arXiv searches on 13 August 2026 included the
exact question, ``cyclic Kaplansky'', weak-* cyclicity, and later papers citing
Lemma 3.3.  No explicit answer or this counterexample was found.  The argument
is elementary, so novelty confidence is moderate rather than high.
\end{itemize}

\paragraph{Human-review recommendation.}
Verify that the intended quantifiers in Question 9.1 are those of Lemma 3.3,
not only the special target $\varphi=1$ mentioned in the preceding prose.
Under the literal wording, the mathematical counterexample is complete.

\begin{thebibliography}{9}
\bibitem{AHMR}
A. Aleman, M. Hartz, J. E. McCarthy, and S. Richter,
\emph{The common range of co-analytic Toeplitz operators on the
Drury--Arveson space}, Journal d'Analyse Math\'ematique 150 (2023), 215--247;
arXiv:2105.01110, revised 2026.

\bibitem{Rudin}
W. Rudin, \emph{Real and Complex Analysis}, 3rd ed., McGraw--Hill, 1987.
\end{thebibliography}

\end{document}
